% TOP_PROBLEM: {"id": 6700062, "title": "Scalar Curvature Question [?66]: Probably, these equalities imply thatP is isometric to a Euclidean rectangular solidbut the approximation/smoo", "category_id": 6, "category": {"id": 6, "name": "geometry", "display_name": "Geometry", "description": "Euclidean and non-Euclidean geometry, geometric structures.", "slug": "geometry", "order_index": 6, "created_at": "2026-07-31T15:26:25.671Z"}, "status": "open"}
% FIRST_POSTED: null
% ENTRY_KIND: statement_only
% Imported from: batch16.tex; UnsolvedMath ID 6700062
\documentclass[11pt]{article}
\usepackage[margin=1in]{geometry}
\usepackage{amsmath,amssymb,mathtools}
\usepackage{fontspec,unicode-math}
\setmainfont{Latin Modern Roman}
\setsansfont{Latin Modern Sans}
\setmathfont{Latin Modern Math}
\usepackage{xurl,hyperref}
\hypersetup{colorlinks=true,linkcolor=blue,urlcolor=blue}
\newcommand{\sref}[2]{\hyperlink{source:#1:#2}{[S#2]}}
\newcommand{\cataloguescope}[1]{\paragraph{Mathematical question.}#1}
\begin{document}
% UnsolvedMath ID 6700062; source catalogue code AMR-066-0062
\setcounter{section}{6700061}
\section{Scalar Curvature Question [?66]: Probably, these equalities imply thatP is isometric to a Euclidean rectangular solidbut the approximation/smoo}\label{6700062}
{\small\sffamily UnsolvedMath ID 6700062 · Geometry}
\subsection{Definitions and mathematical statement}

\paragraph{Shared notation.}$\mathbb N=\{1,2,\ldots\}$, and $\mathbb Z,\mathbb Q,\mathbb R,\mathbb C$ denote the integers, rationals, reals and complex numbers. Any other conventions are specified locally.


For a Riemannian domain, mean curvature $H$ of a smooth boundary face is the trace of its second fundamental form with the convention that the boundary of a Euclidean ball has positive mean curvature. Mean convexity means $H\geq0$. Dihedral angles are interior angles between adjacent faces, measured between their tangent spaces along their common codimension-two faces. Let $P=[0,1]^n$ be the $n$-cube equipped with a Riemannian metric $g$, with its usual faces and corners. Write $\operatorname{Sc}_g$ for scalar curvature, the trace of Ricci curvature, $H_g(Q_i)$ for the mean curvature of each codimension-one face, and $\alpha_{ij}(g)$ for the interior dihedral angles along adjoining faces. In Gromov's 2017 formulation, the hypotheses $\operatorname{Sc}_g\geq0$, $H_g(Q_i)\geq0$ and $\alpha_{ij}(g)\leq\pi/2$ yield the equalities
\[\operatorname{Sc}_g=0,\qquad H_g(Q_i)=0,\qquad\alpha_{ij}(g)=\pi/2.\]
A Euclidean rectangular solid is $\prod_{i=1}^n[0,\ell_i]$ with $\ell_i>0$, equipped with its ordinary Euclidean metric.
\paragraph{Statement recorded in the supplied source.}
Do these equalities imply that $(P,g)$ is isometric to a Euclidean rectangular solid? \sref{6700062}{2}
\subsection{Short English statement}
Must a Riemannian cube with zero scalar curvature, minimal faces and right dihedral angles be a Euclidean rectangular solid?
\subsection{Sources}
\begin{description}
\item[\hypertarget{source:6700062:1}{S1}] ULAM.AI and the UnsolvedMath Contributors, UnsolvedMath: A Curated Collection of Open Mathematics Problems, record 6700062 (AMR-066-0062). CC BY 4.0. \url{https://huggingface.co/datasets/ulamai/UnsolvedMath}.
\item[\hypertarget{source:6700062:2}{S2}] Misha Gromov, \emph{A Dozen Problems, Questions and Conjectures about Positive Scalar Curvature}, arXiv:1710.05946v1 (2017), \S7, assertion [*] and following rigidity question, printed p. 21. \url{https://arxiv.org/abs/1710.05946v1}.
\end{description}
\label{end:6700062}
\end{document}
